\documentclass[11pt]{amsart}
\usepackage[margin=1.1in]{geometry}
\usepackage{amsmath,amssymb,amsthm}
\usepackage{booktabs,array,enumitem,microtype}
\usepackage[hidelinks]{hyperref}
\usepackage{cleveref}
\newcommand{\figdir}{../figures}
% Shared colours and TikZ styles for the 3 x n Col papers.
\usepackage{tikz}
\usetikzlibrary{patterns,arrows.meta,positioning,decorations.pathreplacing,calc,fit,shapes.geometric}

\definecolor{colboth}{RGB}{255,255,255}
\definecolor{colblue}{RGB}{176,206,245}
\definecolor{colbluedark}{RGB}{28,86,176}
\definecolor{colwhite}{RGB}{252,214,160}
\definecolor{colwhitedark}{RGB}{214,128,20}
\definecolor{coldead}{RGB}{150,150,150}
\definecolor{case1}{RGB}{252,187,161}
\definecolor{case2}{RGB}{199,233,192}
\definecolor{case3}{RGB}{198,219,239}
\definecolor{case4}{RGB}{218,218,235}
\definecolor{case5}{RGB}{254,230,206}
\definecolor{case6}{RGB}{247,182,210}

\tikzset{
  gridline/.style={draw=black!30, line width=0.3pt},
  boardline/.style={draw=black!80, line width=0.7pt},
  blockline/.style={line width=1.1pt},
  blocklabel/.style={font=\scriptsize},
  boundlabel/.style={font=\tiny},
  seam/.style={draw=red!70!black, line width=1pt, densely dashed},
  bluestone/.style={circle, fill=colbluedark, draw=colbluedark!70!black, line width=0.4pt, inner sep=0pt},
  whitestone/.style={circle, fill=white, draw=black, line width=0.9pt, inner sep=0pt},
  basebox/.style={draw, rounded corners, fill=case2, font=\small, minimum width=1.5cm, minimum height=0.6cm},
  stepbox/.style={draw, rounded corners, fill=case3, font=\small, minimum width=1.5cm, minimum height=0.6cm},
  rootbox/.style={draw, rounded corners, fill=black!6, font=\small, inner sep=5pt},
  leafbox/.style={draw, rounded corners, fill=case3, font=\small, align=center, inner sep=4pt},
  smallbox/.style={draw, rounded corners, fill=case5, font=\scriptsize, align=center, inner sep=3pt},
}

\usepackage{adjustbox}
\newcommand{\colfig}[1]{\begin{adjustbox}{max width=\linewidth}\input{\figdir/gen/#1.tex}\end{adjustbox}}


\newtheorem{theorem}{Theorem}[section]
\newtheorem{lemma}[theorem]{Lemma}
\newtheorem{proposition}[theorem]{Proposition}
\newtheorem{corollary}[theorem]{Corollary}
\theoremstyle{definition}
\newtheorem{definition}[theorem]{Definition}
\newtheorem{remark}[theorem]{Remark}

\newcommand{\Nb}{N}
\newcommand{\Ebar}{\bar E_4}
\newcommand{\Kc}{K_{\mathrm{corner}}}
\newcommand{\Ktwo}{K_{\mathrm{outer2}}}
\newcommand{\Kfour}{K_{\mathrm{outer4}}}
\newcommand{\Tm}{T_{\mathrm{middle1}}}
\newcommand{\lf}{\mathrel{\triangleleft}}

\title{Every empty $3\times n$ Col board is a second-player win}
\author{Gordon Chi}
\date{\today}

\begin{document}
\begin{abstract}
Col is the partisan placement game in which two players alternately colour
vertices of a graph, a player may never colour a vertex adjacent to one of
their own colour, and the player unable to move loses. We prove that for every
$n\ge1$ the empty $3\times n$ grid has game value $0$, so the second player
wins. Even widths follow from a half-turn mirror strategy. Odd widths use a
one-sided comparison principle that bounds a Col position above by a disjoint
sum of small ``gadget'' positions whenever the cut edges carry no White
legality on both sides. Widths $4k+1$ are covered by a periodic tiling.
Widths $4k+3$ follow by strong induction: three base widths, and a six-case
analysis of Blue's first move in which exactly one case recurses, through a
dead separator column, to a strictly narrower board of the same family. The
finite inputs are eleven gadget inequalities and three base widths, each
certified by a replayable response graph. We give the full written argument, a
seam table and a neighbourhood table that make the width-independence
checkable by hand, and a description of the search-free verifier. The whole
theorem, including every finite certificate, is formalised in Lean~4 without
external libraries; the formal proof uses a streamlined window form of the
same induction that needs no dyadic numbers.
\end{abstract}
\maketitle

\section{Introduction}

\emph{Col} was introduced in Conway's \emph{On Numbers and Games}
\cite{ONAG}, where it is attributed to Colin Vout, and it is a standard
example in combinatorial game theory \cite{WW,Siegel,LIP}. Two players, Blue
(Left) and White (Right), alternately place a stone of their own colour on an
empty vertex of a graph; a player may not place a stone orthogonally adjacent
to a stone of the same colour. Under normal play the player who cannot move
loses. \Cref{fig:rules} shows the rule and a position.

\begin{figure}[ht]
\centering
\colfig{rules_forbidden}\hspace{1.5cm}\colfig{rules_legality}
\caption{Left: the only restriction in Col. Right: a position on the
$3\times5$ board after four moves (numbered). Each empty cell is shaded by who
may still play there; see the legend in \cref{fig:legend}.}
\label{fig:rules}
\end{figure}

Empty rectangular boards are natural test cases. Rectangles with an even
side admit a fixed-point-free half-turn and are second-player wins by
mirroring. For odd-by-odd rectangles no such symmetry exists, since the centre
cell is fixed (\cref{fig:odd}), and the problem becomes delicate.

\begin{theorem}\label{thm:main}
For every integer $n\ge1$, the empty $3\times n$ Col board $H_n$ satisfies
$H_n=0$. Equivalently, the second player wins, whoever moves first.
\end{theorem}

The proof is computer-assisted in a precise and limited sense. A finite list of
local game inequalities (\cref{tab:gadgets}) and three base widths are
certified by explicit response graphs that a short program replays without
search. Everything unbounded is proved in the text: the comparison principle,
the mirror strategy, colour symmetry, the width-independence of the tilings,
and the induction. The finite sweep over widths up to $203$ that the verifier
also runs is a regression test and plays no logical role. Finally, the entire
argument is machine-checked in the Lean~4 proof assistant
(\cref{sec:lean}), with the finite certificates re-checked inside Lean.

\begin{figure}[ht]
\centering
\colfig{roadmap}
\caption{Structure of the proof of \cref{thm:main}.}
\label{fig:roadmap}
\end{figure}

\subsection*{Organisation}
\Cref{sec:prelim} fixes conventions. \Cref{sec:comparison} proves the
comparison principle. \Cref{sec:even} treats even widths, \cref{sec:gadgets}
the gadgets and their certificates, \cref{sec:4k1} widths $4k+1$, and
\cref{sec:4k3} widths $4k+3$. \Cref{sec:locality} proves that the finitely
many local checks cover every width. \Cref{sec:verify} describes the
certificate verifier, \cref{sec:lean} the Lean formalisation, and
\cref{sec:discussion} discusses open problems.

\section{Preliminaries}\label{sec:prelim}

\subsection{Legality sets}
We index cells of the $3\times n$ board by $(r,c)$ with row $r\in\{0,1,2\}$
(top to bottom) and column $c\in\{0,\dots,n-1\}$. For a cell $v$ let $\Nb[v]$
be $v$ together with its orthogonal neighbours. Since the rule depends on a
cell's colour only through legality, a Col position on a graph is determined,
for game purposes, by the pair $(A,B)$ of sets of cells legal for Blue and for
White. Moves act as follows:
\[
\text{Blue at }v\in A:\ (A,B)\mapsto (A\setminus \Nb[v],\ B\setminus\{v\}),\qquad
\text{White at }w\in B:\ (A,B)\mapsto (A\setminus\{w\},\ B\setminus \Nb[w]).
\]
Any pair $(A,B)$ of vertex sets defines a game by these rules, and we call such
a pair a \emph{position} even if no sequence of stones produces it. Gadgets are
positions of this kind. We draw positions as \emph{permission diagrams}
(\cref{fig:legend}): each cell is legal for both players, Blue only, White
only, or neither. These are permissions, not stones.

\begin{figure}[ht]
\centering
\colfig{legend}
\caption{Legend for all board figures. Blue is the Left (positive) player.}
\label{fig:legend}
\end{figure}

\subsection{Game conventions}
We use the standard theory of short partisan games \cite{ONAG,Siegel}. Blue is
Left and White is Right. For games $G,H$: $G+H$ is the disjoint sum, $-G$
swaps the roles, and $G\le H$ means that Right, moving second, wins $H-G$.
In particular $G\le0$ means Blue loses moving first; $G<0$ means White wins
whoever moves first; $G=0$ means the second player wins. We write
$G\lf 0$ for ``not $0\le G$'', that is, White moving first wins. Recall that
$G\le 0$ iff every Left option $G^L$ satisfies $G^L\lf 0$, and $G^L\lf0$ holds
if $G^L$ has a Right option $G^{LR}\le0$ or if $G^L<0$.
Dyadic rationals are games in the usual way, and $1=\{0\mid\ \}$,
$\tfrac14=\{0\mid\tfrac12\}$.

A Col position on a disjoint union of graphs is the disjoint sum of its
restrictions. Swapping every Blue permission with every White permission sends
the position $(A,B)$ to $(B,A)$, which is its negative. The empty board has
$A=B$ and so equals its own negative:

\begin{lemma}[colour symmetry]\label{lem:selfconj}
$H_n=-H_n$. Hence $H_n\le 0$ implies $H_n=0$.
\end{lemma}

\begin{figure}[ht]
\centering
\colfig{game_tree}
\caption{The complete Blue-first game tree of the $1\times 3$ strip, up to
mirror symmetry. Every Blue opening has a White reply after which Blue is
stuck. The strip is also a second-player win when White moves first, so it
has value $0$; this is the case $n=1$ of \cref{thm:main}.}
\label{fig:tree}
\end{figure}

\section{The comparison principle}\label{sec:comparison}

The central tool compares a position with a sum of smaller positions obtained
by cutting the board, provided the cut cannot help White.

\begin{definition}
Let $\Gamma$ be a graph and $P=(A,B)$ a position on $\Gamma$. A \emph{virtual
decomposition} of $P$ consists of a partition of the vertices into blocks,
with the virtual graph $\Gamma^*\subseteq\Gamma$ obtained by deleting every
edge between different blocks, and a position $T=(A^*,B^*)$ on $\Gamma^*$,
such that
\begin{enumerate}[label=(C\arabic*)]
\item $A\subseteq A^*$ (Blue keeps every move),
\item $B^*\subseteq B$ (White gains no move), and
\item no deleted edge has both endpoints in $B^*$.
\end{enumerate}
We write $T=\sum_i T_i$ for the resulting disjoint sum of block positions.
\end{definition}

\begin{theorem}[comparison principle]\label{thm:comparison}
If $T$ is a virtual decomposition of $P$, then $P\le T$.
\end{theorem}

\begin{proof}
We show that Blue, moving first, loses $P+(-T)$. In $-T$ the players'
roles are exchanged: a Blue move of $-T$ is a White placement in $T$, and a
White move of $-T$ is a Blue placement in $T$. White keeps the invariant that
the current positions $P=(A,B)$ and $T=(A^*,B^*)$ satisfy (C1)--(C3).

If Blue plays $v\in A$ in $P$, then $v\in A^*$ by (C1), and White answers in
$-T$ with the Blue placement at $v$ in $T$. Afterwards
$A\setminus\Nb_\Gamma[v]\subseteq A^*\setminus\Nb_{\Gamma^*}[v]$ because
$\Nb_{\Gamma^*}[v]\subseteq\Nb_\Gamma[v]$, and
$B^*\setminus\{v\}\subseteq B\setminus\{v\}$. Condition (C3) persists because
$B^*$ only shrank.

If Blue plays in $-T$, this is a White placement at some $w\in B^*$, and
White answers by playing $w$ in $P$, which is legal by (C2). Then
$A\setminus\{w\}\subseteq A^*\setminus\{w\}$. For (C2), let
$u\in B^*\setminus\Nb_{\Gamma^*}[w]$. Then $u\in B$, and $u\notin\Nb_\Gamma[w]$:
otherwise $uw$ is a deleted edge with both endpoints in $B^*$, contradicting
(C3). Hence $u\in B\setminus\Nb_\Gamma[w]$. Condition (C3) persists.

White therefore always has a reply. The game is finite, so Blue eventually
cannot move and loses. Hence $P-T\le0$.
\end{proof}

\begin{corollary}\label{cor:bound}
Let $T=\sum_iT_i$ be a virtual decomposition of $P$ and suppose
$T_i\le b_i$ for dyadic rationals $b_i$. If $\sum b_i\le 0$ then Blue, moving
first, loses $P$. If $\sum b_i<0$ then $P<0$, so White also wins moving first.
\end{corollary}

\begin{proof}
By \cref{thm:comparison} and monotonicity of addition,
$P\le T\le\sum b_i$.
\end{proof}

\begin{figure}[ht]
\centering
\colfig{sum_intuition}
\caption{Why sums of nonpositive games are nonpositive: White replies in the
component Blue just played in.}
\label{fig:sum}
\end{figure}

\begin{figure}[ht]
\centering
\colfig{comparison_schematic}
\caption{A virtual decomposition, from Case 3 of \cref{sec:4k3} on $3\times15$.
Top: the actual position after Blue's opening and White's reply. Bottom: the
three virtual blocks. Each block may give Blue extra permissions and White
fewer, and at each red seam at most one side is White-legal in every row.}
\label{fig:comparison}
\end{figure}

\begin{remark}
\Cref{thm:comparison} is one-sided. Cutting edges does not in general preserve
value, and (C3) is essential: an edge with White-legal endpoints on both sides
is a restriction on White that the virtual game would forget.
\end{remark}

\subsection{Seams and ports}
In all our decompositions the blocks are intervals of consecutive columns, so
only horizontal edges between the last column of one block and the first
column of the next are deleted. For a block position, its \emph{left and right
White ports} are the White-legality patterns of its first and last columns,
written top to bottom as binary words. Condition (C3) at a seam says exactly
that the right port of the left block and the left port of the right block
have disjoint supports (\cref{fig:seams}).

\begin{figure}[ht]
\centering
\colfig{seams}
\caption{A compatible seam (left) and an incompatible one (right).}
\label{fig:seams}
\end{figure}

\section{Even widths}\label{sec:even}

\begin{proposition}\label{prop:even}
If $n$ is even then $H_n=0$.
\end{proposition}

\begin{proof}
The half-turn $\sigma(r,c)=(2-r,\,n-1-c)$ is a graph automorphism of the
$3\times n$ grid with no fixed cell, since a fixed cell would need
$c=(n-1)/2$. Let the second player answer every move at $v$ by the move at
$\sigma(v)$. We claim that after each answer the position satisfies
$B=\sigma(A)$ (for $A,B$ the Blue and White legal sets, when Blue moved first;
symmetrically otherwise). This holds initially. If Blue plays $v\in A$ then
$\sigma(v)\in B$, and $\sigma(v)\ne v$, so $\sigma(v)$ is still White-legal
after Blue's move. After both moves the new sets are
$A'=A\setminus(\Nb[v]\cup\{\sigma v\})$ and
$B'=B\setminus(\{v\}\cup\Nb[\sigma v])$, and $\sigma(A')=B'$ because $\sigma$
maps $\Nb[v]$ to $\Nb[\sigma v]$ and is an involution. The second player
always has a reply and so wins. Note that $v$ and $\sigma(v)$ may be adjacent;
this is harmless because they receive different colours.
\end{proof}

\begin{figure}[ht]
\centering
\colfig{mirror_even}\hspace{1cm}\colfig{odd_center}
\caption{Left: the half-turn strategy on $3\times 6$; stones $2,4,6$ are
White's rotated answers to $1,3,5$. Right: on odd widths the centre cell is
fixed, and the strategy breaks down.}
\label{fig:odd}
\end{figure}

\section{Gadgets and their certificates}\label{sec:gadgets}

\subsection{The gadgets}
A \emph{gadget} is a position on a $3\times w$ board, given by Blue and White
legality masks. We encode a mask by the integer $\sum 2^{rw+c}$ over its legal
cells. \Cref{tab:gadgets} lists the gadgets and bounds used;
\cref{fig:gallery,fig:opened} draw them. Gadgets with a local move are the
local positions after Blue's opening and White's reply, with White possibly
forgoing further permissions.

\begin{table}[ht]
\centering\small
\begin{tabular}{@{}l c r r c c l@{}}
\toprule
Gadget & $w$ & Blue mask & White mask & ports L/R & bound & used in\\
\midrule
$E_4$ & 4 & 4095 & 2023 & 101/010 & $\le0$ & $4k+1$; Cases 1, 3, 4, 5\\
$\Ebar$ & 4 & 4095 & 3710 & 010/101 & $\le0$ & Cases 2, 3, 4, 6\\
$F_1$ & 1 & 7 & 5 & 101/101 & $\le0$ & $4k+1$\\
$\Kc$ & 3 & 244 & 94 & 011/100 & $\le0$ & Case 2\\
$\Ktwo$ & 7 & 2080241 & 2039675 & 100/101 & $\le0$ & Case 3\\
$\Kfour$ & 7 & 1046471 & 515951 & 101/100 & $\le0$ & Case 4\\
$\Tm$ & 6 & 253500 & 130844 & 001/010 & $\le0$ & Case 5\\
$J_5$ & 5 & 32767 & 21845 & 101/101 & $\le0$ & Case 5\\
$Z_7$ & 7 & 2097022 & 2088828 & 001/101 & $\le0$ & Case 6\\
$C_3$ & 3 & 511 & 503 & 101/111 & $\le\tfrac14$ & Case 1\\
$E_4^{(r,c)}$, $r+c$ odd & 4 & \multicolumn{2}{c}{$4095\setminus\Nb[v]$,\ $2023\setminus\{v\}$} & $\subseteq$101/010 & $\le-1$ & Case 1\\
\bottomrule
\end{tabular}
\caption{The gadget inequalities. $\Kfour$ is the left--right reflection of
$\Ktwo$. The row-2 opened $E_4$ bounds are the top--bottom reflections of the
row-0 ones. A second version of $Z_7$ (White mask 2097020, right port 111)
appears only in the stored base assemblies, at a physical board edge.}
\label{tab:gadgets}
\end{table}

\begin{figure}[ht]
\centering
\colfig{gadget_gallery}
\caption{The gadgets, with their certified bounds. Dots beside each gadget
show its White ports (orange: White-legal). Stones mark the local moves already
played. For $Z_7$ Blue's stone lies in the column just to its left.}
\label{fig:gallery}
\end{figure}

\begin{figure}[ht]
\centering
\colfig{e4_parity}\\[4pt]
\colfig{opened_e4}
\caption{Top: the six odd-parity cells of $E_4$ (dots), those with $r+c$
odd. Bottom: after Blue opens any of them, the remaining position is at most
$-1$.}
\label{fig:opened}
\end{figure}

\subsection{Certificates}\label{sec:certs}
A claim $G+q\le0$, with $G$ a gadget and $q$ a dyadic rational, is certified by
a \emph{response graph}: a finite set $\mathcal S$ of positions of the form
$X+q'$, containing $G+q$, together with a map that assigns to every Left option
of every member of $\mathcal S$ a Right option of that option lying in
$\mathcal S$ (or a terminal position in which Left has no move). By induction on
the finite game, every member of $\mathcal S$ is $\le0$. Checking a response
graph requires only generating Left options and testing membership; no search
is performed. The bound $C_3\le\tfrac14$ is certified as $C_3-\tfrac14\le0$
and $E_4^{(r,c)}\le-1$ as $E_4^{(r,c)}+1\le0$.

\section{Widths congruent to 1 modulo 4}\label{sec:4k1}

\begin{proposition}\label{prop:4k1}
$H_{4k+1}=0$ for every $k\ge0$.
\end{proposition}

\begin{proof}
Decompose the empty board into $k$ copies of $E_4$ followed by $F_1$
(\cref{fig:4k1}). (C1) holds because every block is Blue-legal everywhere,
and (C2) trivially. The seams are $E_4|E_4$ ($010$ against $101$) and
$E_4|F_1$ ($010$ against $101$), so (C3) holds. By \cref{cor:bound},
$H_{4k+1}\le0$, and \cref{lem:selfconj} gives equality.
\end{proof}

\begin{figure}[ht]
\centering
\colfig{tiling_4k1}
\caption{The tiling $E_4^kF_1$, here for $n=13$.}
\label{fig:4k1}
\end{figure}

\section{Widths congruent to 3 modulo 4}\label{sec:4k3}

\subsection{Base cases}
The boards $H_3$, $H_7$, $H_{11}$ are certified directly. For each of their
$4$, $8$ and $12$ openings up to symmetry, the data give a White reply and a
virtual decomposition by certified tiles (the gadgets above and nine further
ordinary tiles), checked by the same comparison conditions.
\Cref{fig:base3} shows the replies on $3\times3$. The six-case step below
does not apply to these widths: at $n=7$ it would fail at $(1,1)$ and
$(1,3)$, and at $n=11$ at $(1,5)$.

\begin{figure}[ht]
\centering
\colfig{base3}
\caption{The four openings of $3\times 3$ up to symmetry and White's certified
replies.}
\label{fig:base3}
\end{figure}

\subsection{The inductive step}
Let $n=4k+3$ with $k\ge3$, and assume $H_m=0$ for every $m<n$ with
$m\equiv3\pmod4$. By \cref{lem:selfconj} it suffices to show $H_n\le0$,
that is, $H_n^L\lf0$ for every Blue opening $H_n^L$.

\subsubsection*{Normalisation}
The top--bottom and left--right reflections are automorphisms of the board
that preserve checkerboard parity $r+c \bmod 2$ (the latter because
$n-1=4k+2$ is even). We may therefore assume that Blue opens at $(r,c)$ with
$r\in\{0,1\}$ and $0\le c\le 2k+1$ (\cref{fig:norm}). Every such opening
falls in exactly one of six cases (\cref{fig:casemap}): odd $r+c$ (Case 1);
$r=0$ and $c=0$ (Case 2), $c\equiv2$ (Case 3), or $c\equiv0$, $c\ge4$
(Case 4) modulo 4; $r=1$ and $c\equiv1$ (Case 5) or $c\equiv3$ (Case 6)
modulo 4.

\begin{figure}[ht]
\centering
\colfig{normalization}
\caption{Normalisation on $3\times15$: by the two reflections every opening
is equivalent to one in the shaded region.}
\label{fig:norm}
\end{figure}

\begin{figure}[ht]
\centering
\colfig{case_map}
\caption{The case used for each normalised opening on $3\times23$.}
\label{fig:casemap}
\end{figure}

In each case we give a decomposition of the position and apply
\cref{cor:bound}; \cref{sec:locality} verifies conditions (C1)--(C3) for all
$k$. Figures \ref{fig:case1}--\ref{fig:case6} show each case on $3\times19$
($k=4$). In every figure the top board is the actual position and the bottom
board the virtual decomposition, with seams dashed.

\subsubsection*{Case 1: $r+c$ odd}
Use $E_4^kC_3$ with no White reply. Since $c\le2k+1<4k$, Blue opened block
$j=\lfloor c/4\rfloor$ at the local cell $(r,c-4j)$, which has odd parity
because $4j$ is even. With $r\in\{0,1\}$ this local cell is one of
$(0,1),(0,3),(1,0),(1,2)$, and that block becomes $E_4^{(r,c-4j)}\le-1$. The
other $E_4$ blocks are $\le0$ and $C_3\le\frac14$, so the position is at most
$-1+\frac14=-\frac34<0$ (\cref{fig:numberline}), and White, to move, wins.

\begin{figure}[ht]
\centering
\colfig{example_case1}
\caption{Case 1, Blue at $(0,5)$ on $3\times19$.}
\label{fig:case1}
\end{figure}

\begin{figure}[ht]
\centering
\colfig{number_line}
\caption{The bookkeeping in Case 1.}
\label{fig:numberline}
\end{figure}

\subsubsection*{Case 2: corner $(0,0)$}
White replies at $(2,2)$. Decompose as $\Kc\,\Ebar^k$; the widths sum to
$3+4k=n$.

\subsubsection*{Case 3: $r=0$, $c=4a+2$}
White replies at $(2,c-2)$. Decompose as $E_4^a\,\Ktwo\,\Ebar^b$ with
$b=k-a-1$; the widths sum to $4a+7+4b=n$. From $4a+2\le2k+1$ we get
$b\ge(2k-3)/4>0$, so $b\ge1$.

\subsubsection*{Case 4: $r=0$, $c=4a+4$}
White replies at $(2,c+2)$. Decompose as $E_4^a\,\Kfour\,\Ebar^b$ with
$b=k-a-1$. Here $a\ge0$, and $4a+4\le 2k+1$ gives $b\ge(2k-1)/4>0$, so
$b\ge1$.

\subsubsection*{Case 5: $r=1$, $c=4a+1$}
White replies at $(0,c-1)$. Decompose as $E_4^a\,\Tm\,E_4^b\,J_5$ with
$b=k-a-2$; the widths sum to $4a+6+4b+5=n$. From $c\le 2k+1$,
$a\le\lfloor k/2\rfloor\le k-2$ for $k\ge3$, so $b\ge0$.

In Cases 2--5 every block is $\le 0$, so after White's reply Blue, to move,
loses by \cref{cor:bound}. Hence $H_n^L$ has a Right option $\le0$.

\begin{figure}[ht]
\centering
\colfig{example_case2}\\[6pt]
\colfig{example_case3}
\caption{Case 2, Blue at $(0,0)$; Case 3, Blue at $(0,6)$ ($a=1$, $b=2$).}
\label{fig:case23}
\end{figure}

\begin{figure}[ht]
\centering
\colfig{example_case4}\\[6pt]
\colfig{example_case5}
\caption{Case 4, Blue at $(0,8)$ ($a=1$, $b=2$); Case 5, Blue at $(1,5)$
($a=1$, $b=1$).}
\label{fig:case45}
\end{figure}

\subsubsection*{Case 6: $r=1$, $c\equiv3\pmod 4$ (the recursive case)}
White replies at $(0,c+1)$. Blue's stone at $(1,c)$ makes all three cells of
column $c$ Blue-illegal. White's reply makes $(0,c)$ White-illegal, and White
voluntarily gives up $(2,c)$, which is allowed by (C2). Column $c$ is then a
\emph{dead separator} (\cref{fig:deadzoom}): it carries no permissions in the
virtual position, so it contributes the game $0$, and no deleted edge at it
has a White-legal endpoint in that column.

On the left, restore every Blue permission lost in columns $0,\dots,c-1$,
which is allowed by (C1). The left block is then exactly the empty board
$H_c$. On the right the width is $s=n-c-1\equiv3\pmod4$, with
$s\ge 4k+2-(2k+1)\ge7$, and we decompose it as $Z_7\,\Ebar^{(s-7)/4}$.
Since $3\le c<n$ and $c\equiv3\pmod 4$, either $c\in\{3,7,11\}$ (base cases) or
$c\ge15$ and the induction hypothesis applies. Either way $H_c=0$, and
\[
\text{position after White's reply}\ \le\ H_c+0+Z_7+\sum\Ebar\ \le\ 0 .
\]

\begin{figure}[ht]
\centering
\colfig{case6_zoom}
\caption{The dead separator column of Case 6.}
\label{fig:deadzoom}
\end{figure}

\begin{figure}[ht]
\centering
\colfig{example_case6}\\[6pt]
\colfig{example_case6_big}
\caption{Case 6. Top: Blue at $(1,7)$ on $3\times19$ recurses to $H_7$.
Bottom: Blue at $(1,11)$ on $3\times23$ recurses to $H_{11}$.}
\label{fig:case6}
\end{figure}

\subsubsection*{Conclusion of the step}
Every normalised opening is covered. In Cases 2--6 White's stated reply leads
to a position $\le0$. In Case 1 the position after Blue's move is already
negative. Thus $H_n\le0$, and $H_n=0$ by \cref{lem:selfconj}. Together with
\cref{prop:even,prop:4k1}, this proves \cref{thm:main}.

\begin{figure}[ht]
\centering
\colfig{recursion_chain}
\caption{Dependencies among the widths $n\equiv3\pmod4$. Each arrow is a
Case 6 call to a narrower board; only the widths $3,7,11$ are certified
directly.}
\label{fig:recursion}
\end{figure}

\section{Width independence}\label{sec:locality}

The decompositions above depend on $k$, $a$ and $b$. We show that
conditions (C1)--(C3) reduce to finitely many checks independent of these
parameters.

\begin{lemma}[locality]\label{lem:locality}
In each case, (C1)--(C3) hold for all admissible parameters provided that:
\begin{enumerate}[label=(\roman*)]
\item every block not containing the opening or the reply has full Blue mask;
\item the exceptional block's masks satisfy (C1) and (C2) against the actual
local position; and
\item every seam that occurs is compatible, and every cell of $\Nb[w]$ for
the White reply $w$ that lies outside the exceptional block is White-illegal
in the virtual position.
\end{enumerate}
\end{lemma}

\begin{proof}
The actual position differs from the empty board only in the cells
$\Nb[v]\cup\{w\}$, removed from $A$, and $\{v\}\cup \Nb[w]$, removed from $B$
(in Case 1 there is no $w$). For a block other than the exceptional one, (C1)
holds by (i). For (C2), a virtual White cell of such a block is White-legal
in the actual position unless it equals $v$ or lies in $\Nb[w]$. The cell $v$
lies in the exceptional block, and the cells of $\Nb[w]$ outside it are
White-illegal virtually by (iii). The exceptional block is (ii), and (C3) is
the seam part of (iii).
\end{proof}

Condition (i) holds because $E_4$, $\Ebar$, $F_1$, $J_5$ and $C_3$ have full
Blue masks, and the left block of Case 6 is the empty $H_c$. Condition (ii)
is a finite check for each of the gadgets with local moves. It is satisfied by
construction, since each gadget's root is the local position after the moves
with some White permissions removed. \Cref{tab:seams,tab:neigh} list every
seam and every cross-block neighbour of a White reply.

\begin{table}[ht]
\centering\small
\begin{tabular}{@{}l l l@{}}
\toprule
Seam (left $|$ right) & ports & where\\
\midrule
$E_4\,|\,E_4$, $E_4\,|\,F_1$, $E_4\,|\,J_5$, $E_4\,|\,C_3$ & $010\,|\,101$ & $4k+1$; Cases 1, 5\\
$E_4^{(r,c)}\,|\,E_4$ or $C_3$ & $\subseteq010\,|\,101$ & Case 1\\
$E_4\,|\,E_4^{(r,c)}$ & $010\,|\subseteq101$ & Case 1\\
$\Ebar\,|\,\Ebar$ & $101\,|\,010$ & Cases 2, 3, 4, 6\\
$\Kc\,|\,\Ebar$ & $100\,|\,010$ & Case 2\\
$E_4\,|\,\Ktwo$,\quad $\Ktwo\,|\,\Ebar$ & $010\,|\,100$,\quad $101\,|\,010$ & Case 3\\
$E_4\,|\,\Kfour$,\quad $\Kfour\,|\,\Ebar$ & $010\,|\,101$,\quad $100\,|\,010$ & Case 4\\
$E_4\,|\,\Tm$,\quad $\Tm\,|\,E_4$ or $J_5$ & $010\,|\,001$,\quad $010\,|\,101$ & Case 5\\
$H_c\,|\,$dead$\,|\,Z_7$,\quad $Z_7\,|\,\Ebar$ & $\cdot\,|\,000\,|\,\cdot$,\quad $101\,|\,010$ & Case 6\\
\bottomrule
\end{tabular}
\caption{Every seam in the six cases and the $4k+1$ tiling. In each row the
two ports have disjoint supports.}
\label{tab:seams}
\end{table}

\begin{table}[ht]
\centering\small
\begin{tabular}{@{}l l l l@{}}
\toprule
Case & reply $w$ (local) & neighbour outside block & White-illegal because\\
\midrule
2 & $(2,2)$ in $\Kc$ & $(2,3)$, first column of $\Ebar$ & left port $010$\\
3 & $(2,0)$ in $\Ktwo$ & $(2,4a-1)$, last column of $E_4$ ($a\ge1$) & right port $010$\\
4 & $(2,6)$ in $\Kfour$ & $(2,4a+7)$, first column of $\Ebar$ & left port $010$\\
5 & $(0,0)$ in $\Tm$ & $(0,4a-1)$, last column of $E_4$ ($a\ge1$) & right port $010$\\
6 & $(0,0)$ in $Z_7$ & $(0,c)$, dead column & dead\\
\bottomrule
\end{tabular}
\caption{Neighbours of White's reply lying outside the exceptional block.}
\label{tab:neigh}
\end{table}

\section{Verification}\label{sec:verify}

The package in \texttt{proofs/construction/three\_row} of the accompanying
repository contains the certificates and a single standard-library Python
verifier, \texttt{verify.py}. It performs no game search and runs in a few
seconds. It checks:
\begin{enumerate}[label=(\arabic*)]
\item SHA-256 hashes of the 34 certificate files;
\item a search-free replay of 21 ordinary response graphs ($20{,}575$
checkpoints, $82{,}899$ edges) and 12 dyadic ones ($1{,}064$ checkpoints,
$3{,}130$ edges), including reflections rebuilt from their sources;
\item the binding of each of the 17 gadget roles to its exact root, ports and
bound;
\item the stored base assemblies for widths $3$, $7$, $11$;
\item as a regression test, all $6{,}651$ assemblies for widths up to $203$,
including $5{,}280$ normalised openings of the six-case step;
\item the local windows (13 plus 2 supplementary row-2 windows) and the 25
seam types, all of which already occur by $n=19$.
\end{enumerate}
A separate script applies 15 targeted corruptions, such as a wrong reply, a
flipped port bit, or a weakened $C_3$ bound. The verifier rejects every one
for the expected reason.

\section{Formal verification in Lean}\label{sec:lean}

The theorem is formalised in Lean~4, using only Lean's core library. The
development is in the directory \texttt{proofs/lean} of the accompanying
repository. The final
statement is phrased with the literal rules of Col:
\begin{quote}\small
\texttt{theorem three\_by\_n\_stones (n : Nat) :}\\
\texttt{\ \ StoneLoses (board n) (fun \_ => False) (fun \_ => False)}
\end{quote}
Here \texttt{StoneLoses R M O} means that the player to move, owning the
stones \texttt{M}, loses against the owner of the stones \texttt{O}, where a
placement is legal if the cell of \texttt{R} is empty and not orthogonally
adjacent to a stone of the mover's colour. Starting from no stones, it states
that the first player loses on every empty $3\times n$ board, whichever colour
starts. The development consists of five modules.

\begin{description}[font=\normalfont\ttfamily, leftmargin=1.5em, style=nextline]
\item[Game] The legality-set model and an inductive predicate
\texttt{Loses}, whose proofs are finite winning strategies for the second
player. It proves the comparison principle (\cref{thm:comparison}, for an
arbitrary subgraph and cut condition), invariance under relabelling of cells,
and the sum rule for non-adjacent regions.
\item[Cert] An encoding of $3\times w$ positions as natural numbers, a
search-free checker for closed certificate sets, and its soundness theorem:
every position in an accepted set is a loss for the player to move.
\item[Blocks] Seams between blocks, repeated $E_4$ and $\bar E_4$ bricks, and
a general \emph{window step} that turns a local domination check around
Blue's opening and White's reply into a proof of the whole position. It is the
locality lemma (\cref{lem:locality}) in general form.
\item[Main] The mirror strategy, the reflections used for normalisation, the
$E_4^kF_1$ tiling, the six cases, and the strong induction.
\item[Stones] The equivalence between the stone rules and the legality-set
model.
\end{description}

\subsection*{A streamlined induction}
The formal proof follows \cref{sec:4k3} but handles all six cases in one
shape (\cref{fig:leanscheme}): the position after Blue's opening and White's
reply is compared with $E_4^a\,W\,\bar E_4^{\,b}$, and in Case 6 with
$H_c\,W\,\bar E_4^{\,b}$. The window $W$ is simply the actual local position
after the two moves. White gives up only the seam cells that the neighbouring
ports require: the middle cell of each edge column next to an $E_4$ or
$\bar E_4$, and the whole first column next to $H_c$ (\cref{fig:leanwindows},
\cref{tab:leanwindows}). Two changes remove all numerical bookkeeping:
\begin{itemize}
\item In Case 1, White replies explicitly at local $(1,4)$ of a
seven-column window starting at the opened brick. The bound
$-1+\tfrac14<0$ is not needed.
\item In Case 5, a seven-column window replaces $\Tm\,E_4^b\,J_5$.
\end{itemize}
With these windows the step already holds for $n=11$ (Case 6 then recurses
to $H_3$), so the only base cases are the whole boards $3\times3$ and
$3\times7$. The even widths and the $4k+1$ tiling are formalised exactly as in
\cref{sec:even,sec:4k1}.

\begin{figure}[ht]
\centering
\colfig{lean_scheme}
\caption{The shape of every case in the Lean proof, shown for Case 3 on
$3\times19$: bricks $E_4^a$, one window containing both moves, and bricks
$\bar E_4^{\,b}$.}
\label{fig:leanscheme}
\end{figure}

\begin{figure}[ht]
\centering
\colfig{lean_windows}
\caption{The nine window gadgets of the Lean proof. Each is the actual local
position after Blue's opening and White's reply, minus the seam cells White
gives up.}
\label{fig:leanwindows}
\end{figure}

\begin{table}[ht]
\centering\small
\begin{tabular}{@{}l c c c c r@{}}
\toprule
Window & width & opening (local) & reply (local) & left neighbour & checkpoints\\
\midrule
Case 1 & 7 & $(0,1),(0,3),(1,0),(1,2)$ & $(1,4)$ & $E_4^a$ & 1456, 1056, 1314, 910\\
Case 2 & 3 & $(0,0)$ & $(2,2)$ & board edge & 7\\
Case 3 & 7 & $(0,2)$ & $(2,0)$ & $E_4^a$ & 1478\\
Case 4 & 7 & $(0,4)$ & $(2,6)$ & $E_4^a$ & 1253\\
Case 5 & 7 & $(1,1)$ & $(0,0)$ & $E_4^a$ & 891\\
Case 6 & 8 & $(1,0)$ & $(0,1)$ & $H_c$ & 4844\\
\midrule
$E_4$, $\bar E_4$, $F_1$ & 4, 4, 1 & & & & 152, 155, 3\\
$H_3$, $H_7$ & 3, 7 & & & & 40, 9819\\
\bottomrule
\end{tabular}
\caption{The certificates checked inside Lean. The right neighbour of every
window is $\bar E_4^{\,b}$.}
\label{tab:leanwindows}
\end{table}

\subsection*{Trusted base}
Everything except the 14 certificate checks is verified by Lean's kernel. The
certificate checks are evaluated by compiled code through Lean's
\texttt{native\_decide}, so the trusted base consists of the Lean kernel, the
Lean compiler, and the axioms \texttt{propext}, \texttt{Classical.choice} and
\texttt{Quot.sound}. The certificates were produced by a separate exhaustive
search, which is not trusted: soundness depends only on the checker. The
build takes about 20 seconds on a laptop.

\section{Discussion}\label{sec:discussion}

The recursive Case 6 is the heart of the argument. A single middle-row stone,
together with a reply and one voluntary concession by White, isolates an
ordinary empty board, which reduces the problem to a narrower width in the
same residue class. We do not know an analogue for five rows. There a single
stone does not kill a column, and experiments suggest that fixed boundary
states are not numerically stable when neutral columns are inserted. The
cases $5\times n$ for odd $n$, and general odd-by-odd rectangles, remain open.

\subsection*{Verification status}
Every part of the proof is machine-checked: the certificates by the Python
verifier of \cref{sec:verify}, and the complete theorem, including its
certificates, by the Lean formalisation of \cref{sec:lean}.

\begin{thebibliography}{9}
\bibitem{WW} E.~R. Berlekamp, J.~H. Conway, R.~K. Guy,
\emph{Winning Ways for Your Mathematical Plays}, 2nd ed., A K Peters, 2001--2004.
\bibitem{ONAG} J.~H. Conway, \emph{On Numbers and Games}, Academic Press,
1976; 2nd ed., A K Peters, 2001.
\bibitem{LIP} M.~H. Albert, R.~J. Nowakowski, D.~Wolfe, \emph{Lessons in
Play: An Introduction to Combinatorial Game Theory}, A K Peters, 2007.
\bibitem{Siegel} A.~N. Siegel, \emph{Combinatorial Game Theory}, Graduate
Studies in Mathematics 146, American Mathematical Society, 2013.
\end{thebibliography}

\end{document}
